\hypertarget{ux6d41ux5339ux914dux6a21ux578b}{%
\chapter{流匹配模型}\label{ux6d41ux5339ux914dux6a21ux578b}}

\hypertarget{ux6d41ux5339ux914dux6a21ux578b-1}{%
\section{流匹配模型}\label{ux6d41ux5339ux914dux6a21ux578b-1}}

\hypertarget{ux4e00ux6d41ux5339ux914dux7684ux57faux672cux601dux8defux53caux4e0eux6269ux6563ux6a21ux578bux7684ux5bf9ux6bd4}{%
\subsection{一、流匹配的基本思路及与扩散模型的对比}\label{ux4e00ux6d41ux5339ux914dux7684ux57faux672cux601dux8defux53caux4e0eux6269ux6563ux6a21ux578bux7684ux5bf9ux6bd4}}

流匹配（Flow Matching）和扩散模型都是生成模型，它们的目标都是将简单的分布（如高斯噪声，\(p_0\)）变换为复杂的数据分布（如图像或文本，\(p_1\)）。

扩散模型的思路是：像墨水在水中扩散一样，数据逐渐加噪变成噪声。生成过程则是逆转这个扩散过程，像``去噪''一样一步步还原数据。这通常被建模为随机微分方程（SDE），路径是离散、曲折的。

流匹配的思路是：希望从加噪后的分布直接得到加噪前的分布。直接定义一个向量场，这个场告诉粒子（数据点）在每一时刻 \(t\) 应该往哪个方向移动。

流匹配中的几个基本概念：

\begin{itemize}
\tightlist
\item
  \textbf{起始状态}：在 \(t=0\) 时刻，所有样本都服从一个简单分布 \(p_0(x)\)，例如标准高斯噪声，这就像一团散沙。
\item
  \textbf{目标状态}：在 \(t=1\) 时刻，样本汇聚成了我们想要生成的复杂数据分布 \(p_1(x)\)，例如精美的图片集合，就像沙子堆成的城堡。
\item
  \textbf{速度场}：记为 \(v(t,x)\)，在任意时刻 \(t\in[0,1]\) 和高维空间的任意位置 \(x\)，存在一个``风向标''或者说``力场''，它指示了位于该位置的粒子在下一瞬间应该往哪个方向移动、速度有多快。
\item
  \textbf{常微分方程的作用}：粒子的运动轨迹 \(x_t\) 随时间的变化率（即速度），正是由该位置的速度场决定的。这在数学上天然地被表达为一个常微分方程：
\end{itemize}

\[
\frac{dx_t}{dt}=v(t,x_t)
\]

如果我们知道每个时刻和位置的风向（即神经网络学到的速度场），我们就必须通过积分（求解这个 ODE）来推导出这粒沙子最终会被吹到哪里。

我们希望让特定数据 \(a\) 从 \(x_0\)（噪声位置，已知）以给定的速度分布（也就是``速度场''中对应位置对应时间的值）趋近 \(x_1\)（有意义的位置，未知）。事实上，上述常微分方程可写成积分形式：

\[
x_1-x_0=\int_0^1 v_t\,dt
\]

由于我们不关心加噪的过程经历了什么，因此我们不妨假设一个理想情况，即所有数据都沿着我们希望的速度和方向匀速直线运动。因此，我们的目标就是要学习这个速度场，使之接近我们假设的理想情况。如果我们把大量噪声点扔进这个场里，它们最终会汇聚成有意义的数据。

\hypertarget{ux4e8cux6570ux5b66ux539fux7406}{%
\subsection{二、数学原理}\label{ux4e8cux6570ux5b66ux539fux7406}}

流匹配的核心在于定义和学习这个``速度场''。

\hypertarget{a.-ux6982ux7387ux8defux5f84probability-path}{%
\subsubsection{A. 概率路径（Probability Path）}\label{a.-ux6982ux7387ux8defux5f84probability-path}}

我们不需要像扩散模型那样模拟复杂的物理扩散过程。我们可以直接定义一条从噪声 \(x_0\) 到数据 \(x_1\) 的路径。最简单且最高效的方法是线性插值（Linear Interpolation）：

\[
x_t=(1-t)x_0+tx_1,\quad t\in[0,1]
\]

\begin{itemize}
\tightlist
\item
  当 \(t=0\) 时，\(x_t=x_0\)，全是噪声。
\item
  当 \(t=1\) 时，\(x_t=x_1\)，全是数据。这条路径是直线的，这对于快速生成非常重要。
\end{itemize}

\hypertarget{b.-ux901fux5ea6ux573avelocity-field}{%
\subsubsection{B. 速度场（Velocity Field）}\label{b.-ux901fux5ea6ux573avelocity-field}}

既然有了位置 \(x_t\) 随时间 \(t\) 的变化公式，我们就可以算出粒子移动的速度（对时间求导）：

\[
u_t(x_t\mid x_0,x_1)=\frac{d}{dt}x_t=x_1-x_0
\]

这个 \(u_t\) 就是条件流匹配（Conditional Flow Matching）的目标速度。它告诉我们：如果要在时间 \(t=1\) 到达 \(x_1\)，当前的粒子应该以什么速度和方向移动。

\hypertarget{c.-ux8badux7ec3ux76eeux6807flow-matching-objective}{%
\subsubsection{C. 训练目标（Flow Matching Objective）}\label{c.-ux8badux7ec3ux76eeux6807flow-matching-objective}}

在训练时，我们无法预知终点 \(x_1\)，这里只是我们想生成的。因此，我们需要训练一个神经网络 \(v_\theta(x_t,t)\)，让它在只知道当前位置 \(x_t\) 和时间 \(t\) 的情况下，去预测上述的目标速度。

损失函数非常简单，就是回归（Regression）问题：

\[
\mathcal{L}_{FM}(\theta)=\mathbb{E}_{t,x_0,x_1}\left[\left\|v_\theta(x_t,t)-u_t(x_t\mid x_0,x_1)\right\|^2\right]
\]

也就是说，让神经网络预测的``流速度''尽可能接近真实的``流速度''。

当然，实际上在图像生成等任务中，预测 \(v\) 时已知的信息包括当前位置 \(x_t\)、时间 \(t\) 和条件 \(c\)。

在给定终点的情况下，我们希望速度应为指向终点的常数。或许你会问，那为什么还需要用神经网络拟合呢？这是因为我们在推理的时候并不知道终点是什么，不同的扩散起点（也就对应不同``粒子''）、不同的附加条件都会对应不同的终点和不同的``理想速度''。但在巨大的高维向量场中，实际在执行特定任务时的速度（该位置的场的方向）往往不是直指向终点的匀速，流匹配模型就是尽量去在每个给定的点拟合我们想要的速度方向。

\hypertarget{ux4e09ux8badux7ec3ux65b9ux6cd5ux4ee5ux6d41ux5339ux914dvlaux6a21ux578bux4e3aux4f8b}{%
\subsection{三、训练方法（以流匹配VLA模型为例）}\label{ux4e09ux8badux7ec3ux65b9ux6cd5ux4ee5ux6d41ux5339ux914dvlaux6a21ux578bux4e3aux4f8b}}

以目前最先进的最优传输流匹配（Optimal Transport Flow Matching，类似于 \(\pi_0\) 中的设计）为例。假设基础噪声分布为 \(\pi_0\sim p_0=\mathcal{N}(0,I)\)，真实动作轨迹分布为 \(\pi_1\sim p_1\)。多模态条件为 \(c=\mathrm{Encoder}(o_t,I_t)\)。

模型构建了一条连接噪声与真实动作的直线概率路径：

\[
x_t=(1-t)x_0+tx_1
\]

动作生成器 \(v_\theta\) 的目标是拟合真实的速度场（Velocity Field），其目标函数定义为：

\[
\mathcal{L}_{FM}(\theta)=\mathbb{E}_{x_0,x_1,t}\left[\left\|v_\theta(x_t,t,c)-(x_1-x_0)\right\|^2\right]
\]

这里 \((x_1-x_0)\) 是真实路径的恒定速度。相较于传统的 DDPM 扩散模型，流匹配的 ODE 轨迹更加平滑且笔直，能够在较少的推理步数（甚至几步内）生成高保真的连续动作信号。

在训练时，我们对于每一组 \(x_0\) 和 \(x_1\)，都希望其连线上每个位置的速度尽可能满足从 \(x_0\) 指向 \(x_1\)，且大小在数值上等于二者之差。当然事实上，由于空间中的速度场是连续的，因此只能尽量拟合上述需求。当多条连线同时汇聚于一个点时，该点速度近似等于各连线对应速度的平均值。

工作流：

\begin{enumerate}
\def\labelenumi{\arabic{enumi}.}
\tightlist
\item
  \textbf{特征提取}：VLM 接收当前相机观测 \(o_t\) 和语言 \(l\)，前向传播生成多模态条件向量 \(c\)。
\item
  \textbf{噪声初始化}：从标准正态分布中采样一段长度为 \(H\) 的轨迹高斯噪声 \(x_0\sim\mathcal{N}(0,I)\)。
\item
  \textbf{常微分方程求解（ODE Solving）}：在时间步 \(t\in[0,1]\) 之间，动作生成器 \(v_\theta(x_t,t,c)\) 结合上下文 \(c\) 计算当前状态的梯度。通过数值积分器（如 Euler 法或 Heun 法）逐步更新 \(x_t\)。
\item
  \textbf{动作执行}：将 ODE 终点 \(x_1\) 作为预测的动作块 \(A_t\)。系统采用``模型预测控制（MPC / Receding Horizon Control）''策略，仅执行前 \(k\) 个动作，随后重新获取新的视觉观测，进入下一轮闭环推理。
\end{enumerate}

这里的时间步是连续的，但神经网络只能处理离散的 \(t\) 值输入。故往往采取采样的方法，如每隔 \(0.02\) 的时间间隔进行一次采样，用该时刻网络输出的值对时间间隔积分，让 \(x_t\)``走一小步''。

\hypertarget{ux56dbux5e94ux7528-1}{%
\subsection{四、应用}\label{ux56dbux5e94ux7528-1}}

\begin{enumerate}
\def\labelenumi{\arabic{enumi}.}
\item
  生成高分辨率、高质量的图像，性能可与Stable Diffusion等扩散模型媲美甚至更优，且通常生成速度更快（因为轨迹更直，从而需要的采样步数更少） 。
\item
  用于生成连续的视频帧序列，处理高维度的时空数据。
\item
  用于科学计算与模拟，如模拟物理系统（如流体动力学）、分子构象生成（蛋白质 3D 结构预测）等，因为这些问题本质上就是对连续变量变化（高维空间中点的移动速度随时空的变化）的建模。
\end{enumerate}

\hypertarget{ux53c2ux8003ux6587ux732e-135}{%
\subsection{参考文献}\label{ux53c2ux8003ux6587ux732e-135}}

\vspace{0.25em}
\begin{itemize}
\tightlist
\item
  Lipman, Y., Chen, R. T. Q., Ben-Hamu, H., Nickel, M., \& Le, M. (2023). \href{https://arxiv.org/abs/2210.02747}{Flow Matching for Generative Modeling}. ICLR.
\item
  Liu, X., Gong, C., \& Liu, Q. (2023). \href{https://arxiv.org/abs/2209.03003}{Flow Straight and Fast: Learning to Generate and Transfer Data with Rectified Flow}. ICLR.
\end{itemize}

\clearpage

\hypertarget{ux79bbux6563ux6d41ux5339ux914d}{%
\section{离散流匹配}\label{ux79bbux6563ux6d41ux5339ux914d}}

\hypertarget{ux4e00ux63d0ux51faux80ccux666f-3}{%
\subsection{一、提出背景}\label{ux4e00ux63d0ux51faux80ccux666f-3}}

目前的生成模型（如Stable Diffusion）在连续空间（像素值是连续的实数）非常成功。但文本是离散的（token ID是整数，如105, 2099），离散空间没有``微小变化''的概念（105变成105.1没有意义），因此直接应用连续扩散模型很困难。

\hypertarget{ux4e8cux6838ux5fc3ux601dux60f3-3}{%
\subsection{二、核心思想}\label{ux4e8cux6838ux5fc3ux601dux60f3-3}}

连续流匹配通过定义一个向量场v\_t(x)，描述样本点x随时间t如何在空间中移动，微分方程为dx/dt=v\_t(x)。在离散空间，我们不能对离散跳变的位置x求导，但概率分布p\_t(x)是随时间连续变化的，我们可以通过连续时间马尔可夫链（CTMC）来建模。

作者发现，如果我们关注从噪声分布p平滑过渡到真实数据分布q的``路径''，定义``概率速度''来描述x的概率分布如何从p转移向q的状态，那么离散版本的公式和连续版本的公式在数学形式上是统一的。这使得很多在连续扩散中好用的技巧（如调度器设计、去噪参数化）可以直接迁移过来。

\hypertarget{ux4e09ux6570ux5b66ux8868ux8fbe-3}{%
\subsection{三、数学表达}\label{ux4e09ux6570ux5b66ux8868ux8fbe-3}}

注意：以下表达式中样本可以取的值的数量（样本空间的元素数）是有限的，概率的值是离散的，p指的是取该样本的实际概率，而不是概率密度函数。

\begin{enumerate}
\def\labelenumi{\arabic{enumi}.}
\tightlist
\item
  概率路径
\end{enumerate}

背景：生成模型的目标是将一个简单的随机噪声分布 \(p_0\)（比如全是掩码 mask 或随机乱码）变成真实的数据分布 \(p_1\)（比如通顺的句子）。

插值 \(p_t\)：我们需要定义一条路径时间 \(t\)（从 \(0\) 到 \(1\)）变化的路径 \(p_t\)：

\begin{itemize}
\tightlist
\item
  当 \(t=0\) 时，\(p_0\) 完全是噪声。
\item
  当 \(t=1\) 时，\(p_1\) 完全是真实数据。
\item
  这条路径是在定义中间过程，在 \((0,1)\) 时，概率分布 \(p_t\) 是如何计算的。
\end{itemize}

\[
p_t(x)=\sum_{x_0,x_1\in\mathcal{D}}p_t(x\mid x_0,x_1)\pi(x_0,x_1)
\]

这个公式描述了整体路径是如何由无数个``个体路径''组成的。
其中，pi(x\_0,x\_1)表示``起始状态为x\_0且终止状态为x\_1的概率''；p\_t(x\textbar x\_0,x\_1)表示``在起始状态为x\_0且终止状态为x\_1的前提下，特定样本x在t时刻出现的概率''（这里x,x\_0,x\_1分别是三个不同样本）。所以这条表达式实际上是全概率公式，通过对不同的(x\_0,x\_1)下在t时刻得到向量x的概率求和，求解出了t时刻采样得到x的总概率。

其中的``个体路径''p\_t(x\textbar x\_0,x\_1)又可以表示为m个概率分布的凸组合：

\[
p_t(x^i\mid x_0,x_1)=\sum_{j=1}^m\kappa_t^j w^j(x^i\mid x_0,x_1)
\]

左边的 \(p_t(x^i\mid x_0,x_1)\) 表示：

\begin{itemize}
\tightlist
\item
  这是目标条件概率。
\item
  它回答了这样一个问题：``如果我们已知起点是 \(x_0\)（比如掩码 \texttt{{[}MASK{]}}），终点是 \(x_1\)（比如单词``猫''），那么在中间某个时刻 \(t\)（比如进度 \(50\%\) 时），第 \(i\) 个位置的 token 取某个值的概率是多少？''
\end{itemize}

右边的 \(\sum_{j=1}^m\) 表示：

\begin{itemize}
\tightlist
\item
  这是一个加权求和。这表明最终的概率分布是由 \(m\) 个基础分布混合而成的。
\end{itemize}

基础分布（Basis Distribution）记为 \(w^j(x^i\mid x_0,x_1)\)，含义如下：

\begin{itemize}
\tightlist
\item
  这是基本的构造块。
\item
  每一个 \(w^j\) 代表一种简单的概率行为。例如：

  \begin{itemize}
  \tightlist
  \item
    \(w^1\) 可能代表``保持为起点 \(x_0^i\)''。
  \item
    \(w^2\) 可能代表``变成终点 \(x_1^i\)''。
  \item
    \(w^3\) 可能代表``变成完全随机的噪声''。
  \end{itemize}
\end{itemize}

因此，该公式定义的是概率的混合：在时刻 \(t\)，这个位置有 \(\kappa_t^1\) 的概率服从分布 \(w^1\)，有 \(\kappa_t^2\) 的概率服从分布 \(w^2\)，以此类推。

对于点x\_i，最终的概率分布是这m个基础分布的加权（凸组合）。若m=2：

\begin{enumerate}
\def\labelenumi{\arabic{enumi}.}
\tightlist
\item
  \(w^1\) 是``保持起点噪声 \(x_0\)''的概率（\(\delta_{x_0}\)）。
\item
  \(w^2\) 是``变成终点数据 \(x_1\)''的概率（\(\delta_{x_1}\)）。
\end{enumerate}

那么公式就变成了：

\[
p_t(x^i\mid x_0,x_1)=(1-\kappa_t)\delta_{x_0^i}+\kappa_t\delta_{x_1^i}
\]

\textbf{物理意义解释}：这句话的意思是，在时间 \(t\)，这个位置的词有 \(\kappa_t\) 的概率已经变成了最终的词 \(x_1\)，有 \(1-\kappa_t\) 的概率还停留在初始的噪声词 \(x_0\)。随着时间 \(t\) 推移，\(\kappa_t\) 逐渐增大，整个序列就慢慢从噪声``流''向了真实数据。

相当于m=2是``在w\_1、w\_2间的线段上取点''，m=3就变成``在w\_1、w\_2、w\_3组成的三角形内取点''，以此类推。几何解释：

假设我们的词表里只有 3 个词：猫、狗、鸟，即 \(d=3\)。

\begin{itemize}
\tightlist
\item
  \textbf{离散状态}：任何一个确定的词（比如``猫''）都可以表示为一个 One-hot 向量：\([1,0,0]\)。
\item
  \textbf{概率分布}：模型输出的不是确定的词，而是一个概率分布，比如 \([0.7,0.2,0.1]\)（\(70\%\) 是猫，\(20\%\) 是狗\ldots\ldots）。
\item
  \textbf{单纯形（Simplex）}：满足``所有元素非负且和为 \(1\)''的向量集合，在几何上构成了一个单纯形。

  \begin{itemize}
  \tightlist
  \item
    对于 3 个词，单纯形是一个三维空间中的等边三角形。
  \item
    三角形的三个顶点分别是 \([1,0,0]\)、\([0,1,0]\)、\([0,0,1]\)，代表确定的词。
  \item
    三角形内部的任何一点都代表一个中间的概率状态。
  \end{itemize}
\end{itemize}

\[
p_t(x^i\mid x_0,x_1)=\sum_{j=1}^{m}\kappa_t^{i,j}w^j(x^i\mid x_0,x_1)
\]

其中的 \(\sum \kappa=1\) 且 \(\kappa\ge 0\)，这正是几何学中``凸组合''的定义。这意味着，无论怎样设计基础路径 \(w^j\)，只要满足这个公式，生成的概率分布 \(p_t\) 永远不会跑出单纯形之外。它保证了模型产生的总是合法的概率分布。

\begin{enumerate}
\def\labelenumi{\arabic{enumi}.}
\setcounter{enumi}{1}
\tightlist
\item
  生成概率速度
\end{enumerate}

\textbf{连续性方程（The Continuity Equation）} 为了定义离散流匹配，我们使用连续时间马尔可夫链（CTMC）范式。样本 \(X_t\) 在离散状态间跳跃。我们需要预测每个 token 概率变化的\textbf{速率}。连续性方程描述了状态概率的变化率 \(\dot{p}_t(x)\) 与通量（flux）之间的关系：

\[
\dot{p}_t(x)+\mathrm{div}_x(p_tu_t)=0
\]

其中 \(\mathrm{div}_x\) 是散度算子。在离散情况下，散度定义为流出通量减去流入通量：

\[
\mathrm{div}_x(v)=\sum_{z\in\mathcal{D}}[v(z,x)-v(x,z)]
\]

\textbf{定理 2（边际速度）} 给定生成条件概率路径 \(p_t(x\mid x_0,x_1)\) 的条件概率速度 \(u_t^i(x^i,z\mid x_0,x_1)\)，其边际速度 \(u_t\) 由下式给出：

\[
u_t^i(x^i,z)=\sum_{x_0,x_1\in\mathcal{D}}u_t^i(x^i,z\mid x_0,x_1)p_t(x_0,x_1\mid z)
\]

这与连续流匹配中的逻辑一致：边际向量场是条件向量场的后验期望。

\textbf{定理 3（闭式解）} 对于由 Eq. 9 定义的路径（即线性插值路径），边际生成概率速度具有非常简洁的闭式解，且与连续流匹配的形式惊人一致：

\[
u_t^i(x^i,z)=\frac{\dot{\kappa}_t}{1-\kappa_t}[p_{1\mid t}(x^i\mid z)-\delta_z(x^i)]
\]

其中 \(p_{1\mid t}(x^i\mid z)=\mathbb{E}[X_1\mid X_t=z]\) 是\textbf{概率去噪器（Probability Denoiser）}。这意味着速度场可以通过学习一个预测目标 \(X_1\) 分布的神经网络来参数化。

\begin{enumerate}
\def\labelenumi{\arabic{enumi}.}
\setcounter{enumi}{2}
\tightlist
\item
  训练目标
\end{enumerate}

\textbf{定理 3 的威力}：它把复杂的``速度计算''转化为了一个简单的监督学习任务。公式 \(u_t\propto p_{1\mid t}-\delta_z\) 告诉我们：只要训练一个神经网络，输入当前带噪声的句子 \(z\) 和时间 \(t\)，预测原始干净句子 \(X_1\) 的概率分布 \(p_{1\mid t}\)，这个网络就自动定义了正确的``流''速度。

\textbf{训练损失函数}：基于上述推导，训练目标简化为最小化交叉熵损失（Cross Entropy）：

\[
\mathcal{L}(\theta)=-\sum_i\mathbb{E}_{t,data}\log p_{1\mid t}(X_1^i\mid X_t;\theta)
\]

这就是标准的\textbf{去噪（Denoising）}目标：给模型看加了噪声/掩码的数据 \(X_t\)，让它预测原数据 \(X_1\)。这解释了为什么 BERT 风格的 Masked Language Modeling（MLM）实际上是在训练一个流模型。

当然，这与BERT的掩码语言模型（MLM）训练有一个关键区别：掩码率是动态的。在t=0时几乎全被掩盖，在t=1时几乎全是真实数据。模型的输入不只有加噪数据X\_t，而是为(X\_t,t)，意味着它需要通过学习来具备在不同噪声水平（与时间t对应）下恢复原始数据X\_1的能力。

\hypertarget{ux53c2ux8003ux6587ux732e-136}{%
\subsection{参考文献}\label{ux53c2ux8003ux6587ux732e-136}}

\vspace{0.25em}
\begin{itemize}
\tightlist
\item
  Gat, I., Remez, T., Shaul, N., Kreuk, F., Chen, R. T. Q., Synnaeve, G., Adi, Y., \& Lipman, Y. (2024). \href{https://arxiv.org/abs/2407.15595}{Discrete Flow Matching}. arXiv:2407.15595.
\item
  Campbell, A., Yim, J., Barzilay, R., Rainforth, T., \& Jaakkola, T. (2024). \href{https://arxiv.org/abs/2402.04997}{Generative Flows on Discrete State-Spaces: Enabling Multimodal Flows with Applications to Protein Co-Design}. arXiv:2402.04997.
\end{itemize}

\clearpage

\hypertarget{ux6d41ux5339ux914dux6a21ux578bux4e2dux7684ux5f3aux5316ux5b66ux4e60ux4e0eflowgrpo}{%
\section{流匹配模型中的强化学习与FlowGRPO}\label{ux6d41ux5339ux914dux6a21ux578bux4e2dux7684ux5f3aux5316ux5b66ux4e60ux4e0eflowgrpo}}

\hypertarget{ux4e00ux9a6cux5c14ux53efux592bux8fc7ux7a0bux7684ux53d8ux5316}{%
\subsection{一、马尔可夫过程的变化}\label{ux4e00ux9a6cux5c14ux53efux592bux8fc7ux7a0bux7684ux53d8ux5316}}

在流匹配模型下，马尔可夫过程的形式和纯自回归的 LLM 有所不同：

\begin{enumerate}
\def\labelenumi{\arabic{enumi}.}
\tightlist
\item
  \textbf{状态 \(x_t\)}：当前时间步的连续图像或潜变量。
\item
  \textbf{动作 \(v_t\)}：模型输出的连续向量场，即一个方向。
\item
  \textbf{环境转移}：确定性的数值积分，例如欧拉法 \(x_{t+\Delta t}=x_t+v_t\Delta t\)。
\item
  \textbf{奖励}：只有积分到 \(t=1\) 得到最终图像 \(x_1\) 后，才能获得美学评分或对齐评分 \(R(x_1)\)。
\end{enumerate}

从 \(x_0\) 到 \(x_1\) 的过程，实质上是 \(v_t\) 关于 \(t\) 从 0 到 1 积分，这里 \(v_t\) 是一个连续函数。事实上一般是采样几十步，每次沿该时间步的速度走一小步。

\hypertarget{ux4e8cux6d41ux5339ux914dux4e2d-rl-ux7684ux7279ux6b8aux96beux70b9}{%
\subsection{二、流匹配中 RL 的特殊难点}\label{ux4e8cux6d41ux5339ux914dux4e2d-rl-ux7684ux7279ux6b8aux96beux70b9}}

前面提到过扩散模型中运用 RL 存在的两个难点：稀疏奖励下的长轨迹信度分配，以及超高维连续动作空间的探索危机。此外，流匹配模型受限于自身底层数学假设，在 RL 中还有一些独有的难题。

\hypertarget{ux7834ux574fux6700ux4f18ux4f20ux8f93ux7684ux76f4ux7ebfux7279ux6027ux4e0eux6d41ux5f62ux8131ux79bbbreaking-optimal-transport-straight-lines}{%
\subsubsection{1. 破坏最优传输的``直线特性''与流形脱离（Breaking Optimal Transport Straight Lines）}\label{ux7834ux574fux6700ux4f18ux4f20ux8f93ux7684ux76f4ux7ebfux7279ux6027ux4e0eux6d41ux5f62ux8131ux79bbbreaking-optimal-transport-straight-lines}}

\begin{itemize}
\tightlist
\item
  \textbf{困境所在}：流匹配，尤其是目前主流的 OT-FM，在预训练时的最大卖点，是通过最优传输（Optimal Transport）强制对齐纯噪声 \(x_1\) 和真实图像 \(x_0\)，构造了一条极其简单的零曲率直线轨迹：
\end{itemize}

\[
x_t=t x_1+(1-t)x_0
\]

\begin{itemize}
\tightlist
\item
  \textbf{直线破裂}：预训练的 FM 是一个完美遵循这条直线的拟合器。当引入 RL 时，RL 策略梯度的唯一目标是``最大化最终图像的奖励''。为了迎合奖励，RL 会强行修改中间时间步的 \(v_t\)，这就不可避免地将 ODE 轨迹``拽离''了原本的最优传输直线。
\item
  \textbf{后果}：扩散模型原本的轨迹就是弯曲的高曲率布朗运动，所以稍微拽偏一点尚能承受。但对于 FM 而言，一旦轨迹脱离直线，不仅会产生偏离流形、生成无意义图像的现象，更致命的是，轨迹曲率的激增会导致 ODE 求解器的截断误差（Truncation Error）爆炸。原本 FM 只需要 4 步积分就能出图，被 RL 扭曲后，必须增加到几十甚至上百步才能维持图像质量，直接粉碎了 FM 最大的``少步数采样优势''。
\end{itemize}

\hypertarget{ux786eux5b9aux6027-ode-ux5e26ux6765ux7684ux4f3cux7136ux5ea6ux8ba1ux7b97ux6b7bux9501likelihood-degeneration}{%
\subsubsection{2. 确定性 ODE 带来的``似然度计算死锁''（Likelihood Degeneration）}\label{ux786eux5b9aux6027-ode-ux5e26ux6765ux7684ux4f3cux7136ux5ea6ux8ba1ux7b97ux6b7bux9501likelihood-degeneration}}

\begin{itemize}
\tightlist
\item
  \textbf{困境所在}：标准在线 RL，如 PPO/GRPO，必须计算动作的对数似然度 \(\log\pi(a\mid s)\)。DM 的逆向去噪本身就是随机高斯转移，天然自带可计算的概率密度。
\item
  \textbf{数学死锁}：FM 的生成过程是纯确定性的常微分方程（ODE）。这意味着给定状态 \(x_t\)，其转移到下一步 \(x_{t-\Delta t}\) 的概率在数学上是一个狄拉克 \(\delta\) 函数（Dirac delta function）。对 \(\delta\) 函数求 log 会导致无穷大（Degenerate Likelihood）。因此，标准 RL 算法在纯粹的流匹配上是完全无法运行的。必须强行魔改底层架构，例如通过 ODE-to-SDE 转换人为注入探索噪声，这让原本极其优雅简洁的 FM 变得异常臃肿和不稳定。
\end{itemize}

当然需要注意的是，流匹配模型在给定输入的情况下输出是固定的，不代表其输出不存在概率分布。在流匹配中，我们常说``去噪过程是固定的''，指的是给定一个具体的初始噪声 \(z\) 后，它演化到目标数据 \(x_1\) 的路径（Trajectory）是唯一且确定的，因为求解的是常微分方程 ODE，而非随机微分方程 SDE。

但是，模型具有分布，是因为输入源头是一个分布。这在数学上被称为前推测度（Push-forward Measure）。可以把标准高斯分布的噪声 \(z\sim\mathcal{N}(0,I)\) 想象成一把沙子，并把沙子撒在一个有特定风向，即向量场 \(v_t(x)\)，的传送带上。虽然每一粒沙子被风吹动的轨迹是完全确定、固定的，但一开始沙子落在各个位置的概率不同。当它们最终被吹到终点时，就会自然而然地堆积成一个新的形状，这就是模型生成的数据分布 \(p_1(x)\)。

\hypertarget{ux4e09flowgrpoux5316-ode-ux4e3a-sde}{%
\subsection{三、FlowGRPO：化 ODE 为 SDE}\label{ux4e09flowgrpoux5316-ode-ux4e3a-sde}}

\hypertarget{ux4e00ode-ux8f6cux5316ux4e3a-sde-ux7684ux65b9ux6cd5}{%
\subsubsection{（一）ODE 转化为 SDE 的方法}\label{ux4e00ode-ux8f6cux5316ux4e3a-sde-ux7684ux65b9ux6cd5}}

\begin{itemize}
\tightlist
\item
  背景原理：在标准的流匹配（Flow Matching）或扩散模型中，图像生成被建模为一个确定性的 ODE。这意味着，只要给定一个初始纯噪声，模型去噪生成图像的轨迹就是唯一确定、不可改变的。
\item
  RL的困境与解法：强化学习（RL）的本质是探索与利用（Exploration and Exploitation）。如果轨迹是确定的，模型就无法尝试新的生成路径，自然也就无从得知哪种路径能获得更高的奖励。因此，FlowGRPO 将这个确定的 ODE 转化为了 SDE。通过在每一步去噪过程中引入由 \(g(t)\) 控制的随机噪声扰动，算法赋予了模型偏离既定路线去``试错''的能力。有了这种随机探索性，RL 才能发挥作用。
\end{itemize}

具体来说，FlowGRPO将流匹配中的常微分方程（ODE）：

\[
dx_t=v_\theta(x_t,t)dt
\]

转化为了随机微分方程：

\[
dx_t=\left[v_\theta(x_t,t)+g(t)\nabla\log p_t(x_t)\right]dt+\sqrt{2g(t)}\,dw_t
\]

仔细观察，可以发现加入的项来自扩散随机加噪的随机梯度朗之万动力学表达式：

\[
dx_t=\nabla\log p(x_t)dt+\sqrt{2}\,dw_t
\]

其中第二项为探索噪声，第一项为保证 \(x\) 关于 \(t\) 的边缘分布不变的修正项，避免模型遇到OOD情形。

乘的系数 \(g(t)\) 表示布朗运动的剧烈程度，是一个预定义好的函数。

\hypertarget{ux4e8cux6bd4ux4f8bux5f52ux4e00ux5316ux4e0eux622aux65adux673aux5236ux7684ux4feeux590d}{%
\subsubsection{（二）比例归一化与截断机制的修复}\label{ux4e8cux6bd4ux4f8bux5f52ux4e00ux5316ux4e0eux622aux65adux673aux5236ux7684ux4feeux590d}}

\begin{itemize}
\tightlist
\item
  背景原理：在 PPO 或 GRPO 等强化学习算法中，为了保证训练的稳定性，必须使用``截断机制''（Clipping）。它通过计算新旧策略的``重要性比值''（Importance Ratio），强制限制模型每次参数更新的幅度，防止模型因为某次偶然的高奖励而彻底破坏原本的生成能力。
\end{itemize}

在流匹配的 SDE 采样中，每一步的去噪转移概率是一个高斯分布。设定当前步的时间间隔为 \(\Delta t\)，噪声方差为 \(\sigma_{t_k}^2\Delta t\)。新策略（当前正在优化的模型）和旧策略（收集数据时的参考模型）的概率密度函数都可以写成高斯形式。重要性比值定义为新旧策略概率之比 \(r_{t_k}(\theta)\)，取自然对数后为：

\[
\begin{aligned}
\log r_{t_k}(\theta)
&=
-\frac{1}{2\sigma_{t_k}^2\Delta t}
\left\|x_{t_k-\Delta t}-\mu_\theta\right\|^2
\\
&\quad+
\frac{1}{2\sigma_{t_k}^2\Delta t}
\left\|x_{t_k-\Delta t}-\mu_{\theta_{old}}\right\|^2
\end{aligned}
\]

关键点来了：在强化学习中，训练数据的轨迹 \(x_{t_k-\Delta t}\) 是由旧策略 \(p_{\theta_{old}}\) 采样生成的。因此它必然满足：

\[
\begin{aligned}
x_{t_k-\Delta t}
&=
\mu_{\theta_{old}}+\sqrt{\sigma_{t_k}^2\Delta t}\cdot\epsilon
\end{aligned}
\]

这里 \(\epsilon\sim\mathcal{N}(0,I)\) 是标准高斯噪声。将这个 \(x_{t_k-\Delta t}\) 代入对数公式，并定义新旧策略的均值差为 \(\Delta\mu=\mu_{\theta_{old}}-\mu_\theta\)，公式展开并化简后得到：

\[
\begin{aligned}
\log r_{t_k}(\theta)
&=
-\frac{\|\Delta\mu\|^2}{2\sigma_{t_k}^2\Delta t}
\\
&\quad-
\frac{\epsilon^T\Delta\mu}{\sigma_{t_k}\sqrt{\Delta t}}
\end{aligned}
\]

对上述公式求高斯噪声 \(\epsilon\) 的数学期望，由于 \(\mathbb{E}[\epsilon]=0\)，第二项被消去：

\[
\begin{aligned}
\mathbb{E}_{\epsilon}\left[\log r_{t_k}(\theta)\right]
&=
-\frac{\|\Delta\mu\|^2}{2\sigma_{t_k}^2\Delta t}
\end{aligned}
\]

因为 \(\|\Delta\mu\|^2>0\)（只要策略更新了，均值就会有差异），所以对数重要性比值的期望严格小于 0。

\begin{itemize}
\item
  截断失效的危机：在流匹配模型中，由于概率密度的特性，这个重要性比值的分布会系统性地向左偏移（平均值小于 1），并且在不同的时间步长上表现出极不一致的方差。这导致预设的截断边界（比如 \([1-\epsilon,1+\epsilon]\)）形同虚设，无法有效约束模型那些过于自信的错误更新，进而引发严重的``奖励黑客''（Reward Hacking）现象，即模型钻空子获得高分但生成一堆无意义的乱象。
\item
  RatioNorm 的解法：为了修复这个问题，论文引入了 RatioNorm 技术。该技术对数重要性比值进行标准化处理，强行将其分布重新居中到零附近。这样一来，截断边界就能再次精准地``卡''住过大的策略更新，确保了最终导出的图像目标函数 \(\mathcal{J}_{Flow}(\theta)\) 能够平稳收敛。
\end{itemize}

\[
\begin{aligned}
\log \tilde{r}_{t_k}(\theta)
&=
\sigma_{t_k}\sqrt{\Delta t}
\left(
\log r_{t_k}(\theta)
\right.
\\
&\quad\left.
+\frac{\|\Delta\mu_\theta(x_{t_k},t_k)\|^2}{2\sigma_{t_k}^2\Delta t}
\right)
\end{aligned}
\]

\hypertarget{ux53c2ux8003ux6587ux732e-137}{%
\subsection{参考文献}\label{ux53c2ux8003ux6587ux732e-137}}

\vspace{0.25em}
\begin{itemize}
\tightlist
\item
  Lipman, Y., Chen, R. T. Q., Ben-Hamu, H., Nickel, M., \& Le, M. (2023). \href{https://arxiv.org/abs/2210.02747}{Flow Matching for Generative Modeling}. ICLR.
\item
  Black, K., Janner, M., Du, Y., Kostrikov, I., \& Levine, S. (2023). \href{https://arxiv.org/abs/2305.13301}{Training Diffusion Models with Reinforcement Learning}. arXiv:2305.13301.
\item
  Liu, J., Liu, G., Liang, J., et al.~(2025). \href{https://arxiv.org/abs/2505.05470}{Flow-GRPO: Training Flow Matching Models via Online RL}. arXiv:2505.05470.
\end{itemize}

\clearpage
